\section{Stake}\label{sec:bond}

The rules of Section~\ref{sec:reward} decide what breadth is worth. They do not decide whose word is taken that
the Bonfires holding a construction are separate parties. That is the work of stake.

\paragraph{Bond and pledge.} The stake its stakeholders hold in a Bonfire is the Bonfire's \emph{bond}.
Stakeholders share in the Bonfire's payment by its own rule, weigh in what enters its grammar, and lose their stake
if the Bonfire is found at fault: they underwrite it, as the masters of a guild answered for its mark. A Bonfire
counts toward the breadth of other Bonfires' constructions only while its bond is at least a minimum, $m$. A node
that serves a Bonfire commits to it a \emph{pledge} of a fixed size, $q$. Stake pledged to one Bonfire cannot be
pledged to another, so where a node places its pledges is a choice with a cost. Because every pledge is the same
size, an operator who runs many nodes has committed many pledges, and the number of nodes behind a Bonfire is a
measure of the stake committed to serving it.

\paragraph{The ceiling.} Stake is the accountability engine of the protocol. It makes a Bonfire and a node
answerable, and it limits what can be earned; it adds nothing to credit. In one epoch a Bonfire is paid at most a
fixed multiple of the pledges committed to it. Writing $k$ for the number of nodes pledged to the Bonfire, $c$ for
the multiple, the \emph{ceiling rate}, and $e$ for what the Bonfire is paid in the epoch,
\begin{equation}\label{eq:ceiling}
  e \;\le\; c\,k\,q .
\end{equation}
Breadth fills the ceiling, and the nodes committed to serving a grammar set it: uptake by other communities is the
demand, and committed capacity is the limit. A pledge earns nothing of itself. The ceiling is linear in the number
of pledges, so dividing or combining nodes changes nothing. A new Bonfire serves itself, and its ceiling is set by its own pledge. Pledges move toward grammars whose breadth earns more than their ceilings admit, and away from those that
cannot fill them, so capital follows uptake and does not direct it. An amount above a ceiling is not paid; what
becomes of it is a question of issuance, left to the token model. The reason for a ceiling and not a weight is
empirical: where stake has raised rewards directly, rewards have followed stake and not service
(Section~\ref{sec:prior}).

A node receives its part of the share $s$, so the most a pledge can yield is $s\,c$ per epoch.

\paragraph{Faults.} Table~\ref{tab:faults} lists what stake answers for.

\begin{table}[!htb]
\centering\small
\begin{tabularx}{\textwidth}{@{}p{0.34\textwidth}p{0.30\textwidth}X@{}}
\toprule
\textbf{Fault} & \textbf{What is lost} & \textbf{How it is decided} \\
\midrule
A receipt that does not re-execute, a receipt withheld on challenge, or a commitment altered after the fact &
The pledge of the node that signed it &
Re-execution by any node; a disagreement is settled by a bonded vote of nodes \\
\addlinespace
One party posing as several Bonfires, so that its own constructions appear widely held &
The bonds of the Bonfires found to be one party, and the pledges of the nodes that party runs &
A judged challenge, settled by a bonded vote \\
\bottomrule
\end{tabularx}
\caption{Faults, what is lost, and how each is decided.}\label{tab:faults}
\end{table}

The first kind of fault is a matter of computation. One honest node suffices to find it, since re-execution has one right answer where the engine is deterministic. A ledger does not run the grammar engine, however, so when the accused node disputes the finding the
dispute goes to a vote of nodes weighted by stake, and the vote is sound while most of the voting stake is honest.
The second kind cannot be settled by computation at all. The question put to the vote is narrow: whether the
governing stake of several Bonfires is controlled by one party. Each real Bonfire has a history to show for
itself, in its members, its own constructions, its gaps, trials and adoptions. Because the protocol is indifferent
to who or what produced a piece of discourse, generated conversation is not a fault. The design of the judged vote
is among the open problems of Section~\ref{sec:open}.

One further rule is proposed and left open. When a Bonfire is found at fault, the nodes pledged to it that the
party at fault does not run lose a part of their pledge. Whoever's capital is exposed is who checks, and a node
that can lose by serving a sham has a reason to examine a Bonfire before pledging to it. The cost is that an
honest node can lose stake through another's fraud, and that nodes will be slow to serve a Bonfire with no history.

A fraction $\varphi$ of slashed stake pays the challenger and the rest is destroyed. If the whole of it went to the
challenger, an owner could challenge its own parties through another identity and lose nothing. Spot checks are
paid from the part of the epoch budget set aside for them, whether or not they find a fault; this, and not the
slash, is what answers the difficulty of Section~\ref{sec:agreements} that unpaid checks are not carried out.

\paragraph{Lock.} A scientist's stake in a claim is reputation, which is slow to build and cannot be lent. A token
can be borrowed for a day. Stake is therefore locked for a period at least as long as a challenge takes to bring,
and whoever has asked to withdraw remains liable for faults committed while the tokens were locked.
